\section{农产品期货数据样本与合约口径}
\label{app:contracts}

本报告涉及的全部期货品种均为交易所挂牌的标准化合约，
主力连续价格序列由新浪财经接口 \texttt{futures\_main\_sina} 获取（日线，取主力合约）。
下表给出各品种的数据样本区间、价格水平与波动特征，
其中“峰谷比”为样本期内历史最高价与最低价之比，
\hl{该指标是判断品种周期振幅的直接依据}。

% 由 scripts/make_tables.py 自动生成，请勿手工编辑
\begin{table}[htbp]
\centering
\scriptsize
\caption{农产品相关期货品种样本概况（主力连续，全部可得历史；末月 2026 年 9 月为不完整月度）}
\label{tab:fut-summary}
\begin{adjustbox}{max width=\textwidth}
\begin{tabular}{@{}lllllrrrrrrrr@{}}
\toprule
品种 & 代码 & 大类 & 起始 & 结束 & 月度数 & 月均价格 & 最新 & 区间涨跌\% & 年化波动\% & 历史最高 & 历史最低 & 峰谷比 \\
\midrule
鸡蛋 & JD & 养殖 & 2013-11 & 2026-09 & 155 & 3\,958.20 & 3\,802.30 & -6.00 & 24.20 & 5\,246.70 & 2\,842.30 & 1.85 \\
生猪 & LH & 养殖 & 2021-01 & 2026-09 & 69 & 16\,205.20 & 11\,381.10 & -55.00 & 27.00 & 28\,171.10 & 10\,420.70 & 2.70 \\
玉米 & C & 粮食 & 2005-01 & 2026-09 & 261 & 2\,081.00 & 2\,239.60 & +93.90 & 11.00 & 2\,975.80 & 1\,154.80 & 2.58 \\
玉米淀粉 & CS & 粮食 & 2014-12 & 2026-09 & 142 & 2\,563.90 & 2\,499.90 & -8.40 & 13.30 & 3\,382.50 & 1\,671.40 & 2.02 \\
粳稻 & JR & 粮食 & 2013-11 & 2022-12 & 110 & 3\,005.70 & 2\,833.00 & -8.10 & 10.50 & 3\,478.00 & 2\,610.90 & 1.33 \\
早籼稻 & RI & 粮食 & 2009-04 & 2022-06 & 159 & 2\,523.00 & 2\,527.50 & +23.90 & 10.80 & 3\,003.60 & 2\,017.20 & 1.49 \\
粳米 & RR & 粮食 & 2019-08 & 2026-09 & 86 & 3\,507.90 & 3\,598.60 & -1.30 & 4.70 & 3\,685.30 & 3\,343.10 & 1.10 \\
强麦 & WH & 粮食 & 2005-01 & 2023-02 & 218 & 2\,464.90 & 3\,230.80 & +92.40 & 11.20 & 3\,443.20 & 1\,554.10 & 2.22 \\
豆一 & A & 油料 & 2005-01 & 2026-09 & 261 & 4\,243.80 & 5\,030.20 & +84.90 & 14.70 & 6\,312.70 & 2\,549.10 & 2.48 \\
豆二 & B & 油料 & 2007-08 & 2026-09 & 230 & 3\,943.30 & 4\,134.20 & +11.50 & 16.00 & 5\,564.70 & 2\,800.40 & 1.99 \\
花生 & PK & 油料 & 2021-02 & 2026-09 & 68 & 9\,008.60 & 7\,951.70 & -24.20 & 12.10 & 10\,882.70 & 7\,798.50 & 1.40 \\
油菜籽 & RS & 油料 & 2012-12 & 2026-09 & 166 & 5\,217.70 & 5\,617.90 & +6.60 & 13.50 & 6\,735.40 & 3\,661.80 & 1.84 \\
豆粕 & M & 压榨链 & 2005-01 & 2026-09 & 261 & 3\,058.50 & 3\,407.00 & +58.40 & 18.30 & 4\,212.90 & 2\,150.30 & 1.96 \\
菜油 & OI & 压榨链 & 2007-06 & 2026-09 & 232 & 8\,542.60 & 10\,223.90 & +18.80 & 16.90 & 13\,595.20 & 5\,568.10 & 2.44 \\
棕榈油 & P & 压榨链 & 2007-10 & 2026-09 & 228 & 6\,971.30 & 10\,071.70 & +18.40 & 22.10 & 11\,687.80 & 4\,251.40 & 2.75 \\
菜粕 & RM & 压榨链 & 2012-12 & 2026-09 & 166 & 2\,529.30 & 2\,349.40 & -0.80 & 19.10 & 3\,949.10 & 1\,869.60 & 2.11 \\
豆油 & Y & 压榨链 & 2006-01 & 2026-09 & 249 & 7\,538.90 & 9\,034.90 & +78.40 & 16.80 & 12\,390.40 & 5\,064.70 & 2.45 \\
棉花 & CF & 软商品 & 2005-01 & 2026-09 & 261 & 15\,853.40 & 16\,335.80 & +27.20 & 16.50 & 32\,782.00 & 10\,271.50 & 3.19 \\
棉纱 & CY & 软商品 & 2017-08 & 2026-09 & 110 & 22\,343.50 & 22\,371.10 & -2.90 & 13.40 & 29\,203.60 & 18\,641.90 & 1.57 \\
20号胶 & NR & 软商品 & 2019-08 & 2026-09 & 86 & 11\,484.20 & 16\,238.30 & +63.70 & 16.10 & 16\,238.30 & 8\,047.10 & 2.02 \\
天然橡胶 & RU & 软商品 & 2005-01 & 2026-09 & 261 & 17\,221.50 & 19\,056.40 & +57.70 & 26.20 & 40\,886.00 & 9\,914.10 & 4.12 \\
白糖 & SR & 软商品 & 2006-01 & 2026-09 & 249 & 5\,382.90 & 5\,400.80 & +3.00 & 13.90 & 7\,344.80 & 3\,032.00 & 2.42 \\
苹果 & AP & 林果 & 2017-12 & 2026-09 & 106 & 8\,252.60 & 7\,321.90 & -9.60 & 23.50 & 11\,975.50 & 5\,540.70 & 2.16 \\
红枣 & CJ & 林果 & 2019-05 & 2026-09 & 89 & 10\,779.30 & 7\,579.70 & -22.00 & 23.30 & 15\,799.10 & 7\,579.70 & 2.08 \\
纸浆 & SP & 其他 & 2018-11 & 2026-09 & 95 & 5\,604.40 & 4\,841.90 & -7.30 & 16.90 & 7\,264.10 & 4\,418.00 & 1.64 \\
\bottomrule
\end{tabular}
\end{adjustbox}
\end{table}


\vspace{10pt}
\noindent\textbf{合约口径说明}：
\begin{enumerate}[leftmargin=1.5em, itemsep=2pt]
  \item \textbf{主力连续合约}由换月拼接而成，\hl{在换月点可能产生价格跳空}，
        因此报告在计算月度收益率时使用月度均价而非月末收盘价，以降低跳空影响。
  \item \textbf{样本端点不一致}：玉米、豆粕、豆油等品种自 \num{2005} 年起有数据，
        而生猪（\num{2021} 年 1 月）、花生（\num{2021} 年 2 月）上市较晚。
        第 \ref{sec:corr} 章的所有相关性计算统一截取
        \num{2021} 年 2 月---\num{2026} 年 8 月，
        以保证跨品种可比。
  \item \textbf{部分品种已停止交易}：强麦（\num{2023} 年 2 月）、
        粳稻（\num{2022} 年 12 月）、早籼稻（\num{2022} 年 6 月）的历史数据
        仍保留在样本中用于长周期分析，但不参与近年相关性检验。
  \item \textbf{纸浆、天然橡胶、棉纱等品种与农业的相关性}：
        纸浆（林产品）、天然橡胶（种植业上游）与棉纱（棉花下游）
        虽不直接属于“农产品”现货范畴，
        但与第 \ref{sec:soft} 章的经济作物种植与加工环节直接相关，
        故纳入样本并单独说明。
\end{enumerate}

\vspace{6pt}
\noindent\textbf{生猪现货数据的补充说明}：
由于生猪期货 \num{2021} 年 1 月才上市，长样本的周期划分
（\num{2015} 年以来）使用生猪现货价格指数（玄田数据/猪易），
并以 \num{20}\% 的锯齿阈值识别周期段；
期货样本的阈值取 \num{18}\%，以匹配其较高的波动率。
两者在重叠区间（\num{2021} 年 2 月---\num{2026} 年 9 月）的周期段划分一致。
